\documentclass[lettersize,journal]{IEEEtran}
\usepackage{amsmath,amsfonts}
\usepackage{algorithmic}
\usepackage{algorithm}
\usepackage{array}
\usepackage[caption=false,font=normalsize,labelfont=sf,textfont=sf]{subfig}
\usepackage{textcomp}
\usepackage{stfloats}
\usepackage{url}
\usepackage{verbatim}
\usepackage{graphicx}
\usepackage{cite}
\usepackage{multirow}
\usepackage{booktabs}
\usepackage{makecell}
\usepackage{tabularx}


\hyphenation{op-tical net-works semi-conduc-tor IEEE-Xplore}
% updated with editorial comments 8/9/2021

\begin{document}

\title{TMFIF-Net: A Two-Stage Multimodal Feature Interaction Fusion Network for Latent Emotion Recognition}

% \author{IEEE Publication Technology,~\IEEEmembership{Staff,~IEEE,}
\author{Shuhuan Zhao\thanks{}, Jiaqi Chen, Shuwen Yang, Ruiyang Duan, Zhen Wang, Shuaiqi Liu
        % <-this % stops a space
\thanks{This work was supported in part by National Natural Science Foundation of China under Grant 62172139, Science and Technology Program of Hebei Province under Grant 246Z0104G, Shijiazhuang Key Research and Development Special Project for Universities in Hebei Province under Grant 2511300801A, Hebei University Research and Innovation Team Support Project under Grant IT2023B05. This work was also supported by the High-Performance Computing Center of Hebei University. (\textit{Corresponding author: Shuaiqi Liu)}}
\thanks{S. Zhao, and S. Liu are with College of Electronic and Information 
Engineering, Hebei University, 071002, China, and also with the Machine Vision Technology Innovation Center of Hebei Province, Baoding 071002, China (e-mail: zsh\_hbu@ 163.com; shdkj-1918@163.com).}
\thanks{J. Chen, S. Yang, and R. Duan are with College of Electronic and Information 
Engineering, Hebei University, 071002, China (e-mail: jiaqi\_chen2024@163.com; yswsxka@163.com; 2964521764@qq.com).}
\thanks{Z. Wang is with Hebei Xiong'an Taoli Technology Co., Ltd, Baoding 071002, China (e-mail: hbu\_ee@163.com).}


}% <-this % stops a space
%\thanks{Manuscript received April 19, 2021; revised August 16, 2021.}%


% The paper headers
%\markboth{Journal of \LaTeX\ Class Files,~Vol.~14, No.~8, August~2021}%
\markboth{IEEE Transactions on Affective Computing, 2026}
{ Zhao \MakeLowercase{\textit{et al.}}: TMFIF-Net: A Two-Stage Multimodal Feature Interaction Fusion Network for Latent Emotion Recognition}




%\IEEEpubid{0000--0000/00\$00.00~\copyright~2021 IEEE}
% Remember, if you use this you must call \IEEEpubidadjcol in the second
% column for its text to clear the IEEEpubid mark.

\maketitle

\begin{abstract}
Micro-expressions are regarded as reliable cues of genuine emotions, which have significant applications in fields such as healthcare and criminal investigations. However, relying on a solitary facial muscle movement in real-world may provide insufficient information for accurate recognition. In this paper, we propose a Two-stage Multimodal Feature Interaction Fusion Network (TMFIF-Net) for latent emotion recognition. The strengths of RGB, depth, and ECG signal are integrated, with heterogeneity and correlations across modalities taken into account, and cross-modal complementarity leveraged to improve recognition accuracy. The overall framework consists of unimodal feature extraction and multimodal fusion. To overcome the constraints of unimodal micro-expression recognition, a Hierarchical Spatial Feature Extraction (HSFE) block is devised to extract discriminative shallow features for each modality. Subsequently, a Dual-Branch Feature Aggregation (DBFA) block is established to extract fine-grained local and global context. A two-stage fusion strategy is proposed to mitigate modality heterogeneity. In the first stage, an Adaptive Dynamic Weighting Fusion (ADWF) block is put forward to accomplish the initial cross-modal fusion. In the second stage, a Bidirectional Guided Interaction Fusion (BGIF) block is devised to discern the correlations among modalities and augment the richness of the fused features. The proposed method achieves notable performance improvements on the CAS(ME)³ dataset, with UF1 and UAR of 0.7124 and 0.6633, respectively.
\end{abstract}

\begin{IEEEkeywords}
Affective computing, emotion recognition, micro-expression recognition, multimodal.
\end{IEEEkeywords}

\section{Introduction}
\IEEEPARstart{F}{acial} expressions generally encompass macro-expressions (MAE) and micro-expressions (ME), serving as the most direct modality, which enables humans to perceive emotions. MEs differ from macro-expressions\cite{1} in that they occur involuntarily, exhibit lower intensity, and last only briefly (0.04-0.2 seconds)\cite{2}. Thus, distinguishing them from feigned expressions remains nontrivial due to the potential confusion. Existing studies indicate that high recognition accuracy is rarely achieved, even by experts receiving extensive professional training. Beyond facial expressions, affective states are also naturally conveyed through body movements, voice, depth information, and physiological signals. Therefore, multimodal micro-expression recognition (MER) is essential for the development trend in affective computing. Different modalities reveal emotional changes from distinct perspectives. The depth information can provide the three-dimensional structure changes of the face, which is helpful for capturing the facial contour changes caused by muscle movements\cite{3}. When triggered by emotional arousal, physiological signals become beyond conscious control and generate internal fluctuations, making it difficult to deliberately conceal genuine emotion\cite{4}. The ECG signals provide a reliable measure of heart rate variability, and the emotional arousal and autonomic nervous system regulation are reflected in heart rate changes, serving as one of the ways to convey implicit emotions \cite{5,6}. By integrating visual, physiological signals, and other multimodal information during the generation of emotion, which can enhance the ability to perceive genuine emotion\cite{7}. Applications and research interest extend to public security\cite{8}, human-computer interaction\cite{9}, healthcare\cite{10}, and education\cite{11}. 

Researchers have conducted a series of emotion analyses based on multimodal fusion. These methods can be classified into early fusion\cite{12}, intermediate fusion\cite{13}, and late fusion\cite{14,15}. Early fusion captures subtle correlations among modalities, yet it necessitates modal alignment. In the absence of such alignment, it may introduce irrelevant noise and struggle to handle modalities with notably distinct feature distributions. Intermediate fusion promotes the independence of unimodal feature extraction and the interaction between multiple modalities, but mitigating overfitting generally requires adequate training samples. Late fusion preserves the specificity of each modality's features, but it may lose the semantic information at intermediate levels.

In summary, this paper integrates facial images, depth information, and ECG signals, comprehensively exploiting the heterogeneity and correlation of multimodal data, and proposes a Two-stage Multimodal Feature Interaction Fusion Network (TMFIF-Net) for multimodal latent emotion recognition. The main contributions are as follows: 

(1) To address the limitations of unimodal emotion features in comprehensively revealing complex emotional states, the TMFIF-Net is proposed for multimodal latent emotion recognition. By leveraging the complementary emotional characteristics across visual and physiological dimensions of RGB, Depth, and ECG modalities, TMFIF-Net offers a robust multimodal framework for emotion recognition.

(2) To overcome insufficient intra-modal emotion representation, a Hierarchical Spatial Feature Extraction (HSFE) block is designed to strengthen shallow affective cues, while a Dual-Branch Feature Aggregation (DBFA) block is further constructed to jointly model both local details and global semantics, thereby providing a more comprehensive heterogeneous emotion representation.

(3) To alleviate the limited inter-modal interaction, a two-stage strategy is presented, in which an Adaptive Dynamic Weighting Fusion (ADWF) block assigns modality-specific weights according to modality contribution differences for initial fusion, followed by a Bidirectional Guided Interaction Fusion (BGIF) block that models inter-modal dependencies via attention mechanisms, enabling more precise cross-modal emotion association integration.

\section{Related Works}
\subsection{Unimodal Micro-Expression Recognition}
MER has increasingly been addressed using end-to-end deep neural networks as deep learning progresses, enabling automatic feature extraction with superior performance compared to traditional handcrafted methods. For example, Liu et al.\cite{16} proposed Main Directional Mean Optical-flow (MDMO) and segmented facial regions based on Region of Interest (ROI) using action units (AU) to capture local motion information. Gan et al.\cite{17} developed the Off-ApexNet model, which incorporated a framework with optical flow computed between the onset and apex frames to extract facial expression features. Liong et al.\cite{18} designed a Shallow Triple Stream Three-dimensional CNN Network (STSTNet) to distinguish ME details by leveraging three extracted optical flow features, including horizontal, vertical, and optical strain. Zhou et al.\cite{19} computed the horizontal and vertical optical flow feature between two frames using the TV\_L1 method, which were then fed into a dual-inception network, significantly alleviating cross-domain discrepancies. Xia et al.\cite{20} adopted a shallow Recurrent Convolutional Network (RCN), and subtle dynamic cues latent in low-resolution optical-flow images were captured. This design improved recognition accuracy with reduced computational complexity. Chen et al.\cite{21} introduced the Block Division Convolutional Network (BDCNN) with optical flow features as inputs. Feature vectors were extracted separately from multiple blocks, facilitating facial expression region learning. Cai et al.\cite{22} constructed a dual-branch learning framework of Multi-level Flow-driven Attention Network (MFDAN) for MER, a Transformer-based architecture guided by flow-driven attention. ROI localization is performed to obtain more sensitive and fine-grained ME features. Wang et al.\cite{23} proposed a Hierarchical Transformer Network (HTNet), in which optical flow images were partitioned into four regions, where local self-attention was applied. Cross-region dependencies were modeled through the aggregation layers, which enabled interactions across blocks, enhancing both fine-grained and coarse-grained feature extraction.

MER in affective computing relies on unimodal feature extraction as a basic component. Unimodal approaches have been explored in existing work to learn emotion-related features, yielding improved feature representations. However, local details and global semantic structure are rarely captured simultaneously in most methods, leading to reduced discriminative ability and robustness of emotional features within a modality. Facial expressions are a key carrier for emotional expression and provide important cues to an individual’s psychological state. Considerable progress has been achieved in MER using unimodal RGB data. However, RGB images offer limited affective cues in real-world scenarios. Subtle emotional variations may be missed when only the RGB modality is considered. In contrast, multimodal MER integrates complementary information from multiple modalities. This integration provides a more comprehensive view and improves robustness and representational capacity in complex scenarios. In order to achieve better MER effect, in this paper, we combine the modal advantages of RGB, depth information, and ECG signals, consider the heterogeneity and correlation of modalities, and construct TMFIF-Net.

\subsection{Multimodal Emotion Recognition}

Emotional states are multi-dimensional, complex processes involving psychological and physiological changes. To overcome the limitations of unimodal approaches, multimodal models integrate facial expressions, depth information, and physiological data. Multimodal collaborative interaction makes emotional recognition more reliable. The CAS(ME)3 \cite{24} dataset, released in 2022, contains RGB-D ME images synchronized with physiological recordings, facilitating studies on MER using cross-modal fusion.

At present, scholars from various countries have started research on multimodal MER. For example, Nahid et al.\cite{25} integrated facial, linguistic, and physiological signals for emotion recognition. Their findings highlighted the advantage of multimodal designs for automated sentiment analysis. Xie et al. \cite{26} developed a dialogue emotion framework. Audio, video, and text were fused using Transformer-based cross-modal interactions together with the EmbraceNet architecture. This method showed improved accuracy compared with unimodal approaches. However, the issue of data imbalance was not effectively mitigated. Yin et al.\cite{27} proposed a Token-disentangling Mutual Transformer (TMT) to handle unequal modality contributions. TMT disentangled consistency features across modalities and heterogeneity features within each modality to improve multimodal recognition robustness. However, it failed to resolve the issue of missing emotion labels. Liong et al.\cite{28} proposed the end-to-end attention network based on RGB-D scene flow. This design extended MER to the 3D motion domain. Their results indicated that combining optical flow with depth information helped emphasize salient facial-motion regions and improved emotion prediction performance. Zhang et al.\cite{7} introduced a multimodal learning framework that combines MEs with physiological signals for latent emotion recognition using RGB, depth, and physiological inputs. However, their results indicated that incorporating depth led to lower performance than fusing RGB with physiological signals alone. Zhang et al.\cite{3} constructed a lightweight point cloud graph convolutional network to segment facial regions based on structural and motion cues. Local and global features were captured over the segmented regions, resulting in improved accuracy and robustness for MER. Zhao et al.\cite{29} proposed a quantification framework for multimodal emotion recognition. Multi-scale convolutions and feature-contracting attention were adopted to capture subtle spatiotemporal ME cues. Global dependencies in ECG physiological signals were modeled with Transformer, facilitating the identification of typical composite states. Li et al.\cite{30} integrated ECG and PPG signals with facial expressions to achieve multimodal emotion recognition. Their multimodal fusion approach demonstrated significantly higher accuracy than unimodal methods.

In recent years, multimodal emotion recognition methods have achieved notable progress in integrating information from multiple modalities. However, direct fusion approaches may result in insufficient information interaction. Transformer-based attention mechanisms have been investigated in prior studies to mitigate this limitation, which effectively enrich emotional representations by focusing on distinct modal features, thereby enhancing the consistency of multimodal emotion feature modeling. Nevertheless, data distributions and expression patterns differ markedly across modalities. Across different emotional states, the relative importance of each modality shifts accordingly. Consequently, cross-modal dependencies modeling remains challenging under such conditions. This motivates the development of more efficient mechanisms for cross-modal interaction and fusion. To overcome insufficient intra-modal emotion representation, we construct the HSFE block and DBFA block to provide a more comprehensive heterogeneous emotion representation. We also give a two-stage strategy to alleviate the limited inter-modal interaction.

\section{Methodology}

To mitigate the limitations of unimodal MER, a tri-modal input strategy is adopted to construct the TMFIF-Net, as illustrated in Fig.~\ref{fig_1}. The inputs include optical flow motion cues extracted from ME frame sequences, depth information for spatial guidance, and ECG signals offering physiological changes. The proposed framework is organized into three components: data preprocessing, unimodal feature extraction, and multimodal feature fusion. Unimodal feature extraction is comprised of the HSFE block and the DBFA block, designed to efficiently extract emotion-discriminative features within each modality. In the multimodal fusion, two blocks, including the ADWF and BGIF blocks, are designed to accomplish two-stage fusion. In the first fusion stage, the ADWF block assigns weights to each modality to establish preliminary information integration. In the second fusion stage, the BGIF block introduces an attention mechanism to guide focus on crucial emotion-related information.

\begin{figure*}[!t]
\centering
\includegraphics[width=\textwidth]{fig1}
\caption{The network architecture diagram of the TMFIF-Net, which comprised data preprocessing, unimodal feature extractor, and multimodal feature fusion.}
\label{fig_1}
\end{figure*}

\subsection{Data Preprocessing}

In this paper, a trimodal input strategy consisting of RGB, ECG, and depth information is employed. The preprocessing of RGB images is shown in Fig.~\ref{fig_2}. Since the CSA(ME)$^{3}$\cite{24} dataset only provides the indexes of onset frame $F_{\mathit{c}}^{onset}$ and offset frame $F_{\mathit{c}}^{offset}$ of each expression sequence, without offering index of the apex frame $F_{\mathit{c}}^{apex}$. Thus the frame difference method \cite{31} is utilized to index the apex frame. In this paper, we use MediaPipe\cite{32} method to detect and spot face regions across $F_{\mathit{c}}^{onset}$ and $F_{\mathit{c}}^{offset}$  based on 468 facial landmarks, which is cropped to a size of 224×224 to eliminate background noise information. MEs involve subtle changes in facial muscles, making it difficult to directly extract motion information from $F_{\mathit{c}}^{apex}$. In this paper, we use the TV\_L1\cite{33} method to estimate optical flow motion $F_o = [u, v, s] \in \mathbb{R}^{3 \times 224 \times 224}$  between $F_{\mathit{c}}^{onset}$  and $F_{\mathit{c}}^{apex}$, where $(u, v, s)$ denote horizontal, vertical, and optical strain respectively. 

\begin{figure}[!t]
\centering
\includegraphics[width=2.7in]{fig2}
\caption{The preprocessing of the RGB images.}
\label{fig_2}
\end{figure}

The preprocessing steps for ECG signals involve data alignment, cropping, normalization, and Gramian Angular Difference Fields (GADF)\cite{34} transformation, as illustrated in Fig.~\ref{fig_3}. The corresponding ECG signal segments are synchronized with the RGB ME index calibration in the dataset. Firstly, data alignment is performed on the raw ECG signal corresponding to the video frame rate to acquire $S = [s_1, s_2, \ldots, s_L]$ (Fig.~\ref{fig_3}(a)), where $L$ denotes the length of the alignment reference range. Secondly, the $S$ is cropped according to the indexes from the onset to the offset of the RGB images to obtain $S_e = [s_{\mathit{onset}},\ s_{\mathit{onset}+1},\ \ldots,\ s_{\mathit{offset}}]$ (Fig.~\ref{fig_3}(b)). Subsequently, ECG segments aligned with facial expressions $S_e$  are normalized to acquire $S_{\mathit{norm}} = [s_{\mathit{norm}}^{\mathit{onset}}, s_{\mathit{norm}}^{\mathit{onset}+1}, \ldots, s_{\mathit{norm}}^{i}, \ldots, s_{\mathit{norm}}^{\mathit{offset}}]$ (Fig.~\ref{fig_3}(c)), with values scaled to the range of $-1 \le s_{\mathit{norm}}^{i} \le 1$. The normalization process can be expressed as:

\begin{equation}
\label{eq_1}
\begin{split}
S_{\mathit{norm}}^{i}
=
\frac{s_i - \mathit{min}\!\left(S_e\right)}
     {\mathit{max}\!\left(S_e\right) - \mathit{min}\!\left(S_e\right)}
\times 2 - 1,
\quad \\
\quad i \in \{\mathit{onset}, \mathit{onset}+1, \ldots, \mathit{offset}\}
\end{split}
\end{equation}

\noindent where, $\mathit{max}(\cdot)$ is the function for taking the maximum value of a sequence, while $\mathit{min}(\cdot)$ is the function for taking the minimum value of a sequence.

\begin{figure}[!t]
\centering
\includegraphics[width=2.5in]{fig3}
\caption{The preprocessing of the ECG Signal.}
\label{fig_3}
\end{figure}

Image-based representations derived from time series perform well in classification tasks\cite{35}. The GADF method emphasizes temporal contrast and dynamic variations, which are important for detecting abnormal changes in ECG signals\cite{36}. It also captures local fluctuations triggered by MEs. Therefore, $S_{\mathit{norm}}^{i}$ is converted into angle values bounded in $[0, \pi]$ as follows:

\begin{equation}
\label{eq_2}
\theta_i = \mathit{arccos}\!\left(S_{\mathit{norm}}^{i}\right)
\end{equation}

\noindent where, $\theta_i$ denotes the angle value of the signal mapping.

Finally, based on the GADF method, the corresponding two-dimensional matrix $F_e \in \mathbb{R}^{1 \times 224 \times 224}$ (Fig.~\ref{fig_3}(d)) for the ECG signal is generated as:

\begin{equation}
\label{eq_3}
\small
F_e =
\begin{pmatrix}
\mathit{sin}(\theta_1 - \theta_1) & \mathit{sin}(\theta_1 - \theta_2) & \cdots & \mathit{sin}(\theta_1 - \theta_n) \\
\mathit{sin}(\theta_2 - \theta_1) & \mathit{sin}(\theta_2 - \theta_2) & \cdots & \mathit{sin}(\theta_2 - \theta_n) \\
\vdots & \vdots & \ddots & \vdots \\
\mathit{sin}(\theta_n - \theta_1) & \mathit{sin}(\theta_n - \theta_2) & \cdots & \mathit{sin}(\theta_n - \theta_n)
\end{pmatrix}
\normalsize
\end{equation}




\noindent where, $\mathit{sin}(\theta_i - \theta_i)$ denotes the inherent variability of the time series signal points themselves, taking a constant value of 0. $\mathit{sin}(\theta_i - \theta_j)$ and $\mathit{sin}(\theta_j - \theta_i)$ exhibit an antisymmetric structure, representing the dynamic changes of ECG, where $i, j \in \{1, 2, \ldots, n\}$.

Zhang et al. \cite{7} indicated that, despite combining depth maps with RGB images and physiological signals, directly incorporating depth maps as input did not significantly improve classification performance. In contrast, Li et al.\cite{37} employed frame-difference to process depth maps, finding that this method demonstrated effective synergy with the RGB modality. During the depth images preprocessing illustrated in Fig.~\ref{fig_4}, the frames of depth images corresponding to the onset and apex in each micro-expression sequence are selected and denoted as $F_{\mathit{d}}^{onset}$, $F_{\mathit{d}}^{apex}$. The difference $F_d \in \mathbb{R}^{1 \times 224 \times 224}$ between these two frames is regarded as one input of our network and can be calculated as:

\begin{equation}
\label{eq_4}
F_d = F_{\mathit{d\_apex}} - F_{\mathit{d\_onset}}
\end{equation}

\begin{figure}[!t]
\centering
\includegraphics[width=1.7in]{fig4}
\caption{The preprocessing of the depth images.}
\label{fig_4}
\end{figure}

\subsection{Unimodal Feature Extraction}

Following data preprocessing, the feature maps $F_m,\ m \in \{o, e, d\}$, where $o, e, d$ denote optical flow, ECG signals, and depth data, respectively, which is prepared for subsequent unimodal feature extraction. In this phase, the HSFE block is devised to extract low-level features from each modality. Additionally, the DBFA block utilizes two parallel branches, facilitating intra-modal cross-branch feature complementarity to augment the richness and discriminative capacity of unimodal features. 

\subsubsection{The HSFE Block}


The structure of the HSFE Block is presented in Fig.~\ref{fig_5}, comprises one Two-Conv Layer Block (TCLB) and two Depthwise Convolution Residual Blocks (DCRB). First, a 3×3 convolution layer is employed, after which the size of $F_{\mathit{m}}$ is reduced to half of its original dimensions. Thereafter, it undergoes further processing via LayerNorm and the GELU. Then, the size of $F_{\mathit{m}}$ is decreased by 1/4 to acquire $F_{\mathit{s}}$ via a 3×3 convolution operation and LayerNorm. Subsequently, two DCRBs are connected. In the DCRB, after a 3×3 depthwise convolution, the output serves as the input to the subsequent normalization, a linear layer, and GELU activation function to obtain $F_{\mathit{g}}^{i}$. Simultaneously, a learnable per-channel scaling vector $\lambda_i \in \mathbb{R}^{C}$ is employed to reweight the channel information, resulting in $F_{\lambda}^{i} = F_{g}^{i} \times \lambda_i$, thereby ensuring stable gradient propagation during the initial training phase. $F_{\mathit{\lambda}}^{1}$ and  $F_{\mathit{s}}$ are combined through a residual connection to derive the first-layer features $F_m^{1} = F_{\lambda}^{1} + F_s$. Finally, $F_{\mathit{m}}^{1}$ is fed into the second-layer DCRB, and the above procedure is repeated to obtain the final output features $F_m^{h} = F_{\lambda}^{2} + F_{\mathit{m}}^{1}$. 

\begin{figure}[!t]
\centering
\includegraphics[width=2.5in]{fig5}
\caption{The network architecture diagram of the HSFE block.}
\label{fig_5}
\end{figure}



\subsubsection{The DBFA Block}

The DBFA Block is designed to further extract features of $F_m^{h},\ m \in \{o, e, d\}$, consisting of a global branch and a local branch, as depicted in Fig.~\ref{fig_6}. For the global branch, all three modality features share the identical structure. The global branch is based on the Swin-Transformer\cite{38} hierarchical architecture, consisting of four stages. Each stage encompasses two Lite SwinBlocks (LSBlock) (Fig.~\ref{fig_6}(a)), which employ an alternating scheme of the Window Multi-head Self Attention (W-MSA) and the Shifted Window Multi-head Self Attention (SW-MSA) to capture global contextual information. Stage 1 does not perform downsampling, preserving the shallow feature dimensions to obtain the feature $F_G^{1}$. Stages 2-4 implement hierarchical downsampling operations, reducing the spatial resolution of feature maps to half that of the preceding layer while doubling the number of channels. The final representation obtained from the global branch is denoted as $F_G^{4}$. The local branch consists of a Depthwise Convolution Attention Block (DCABlock) (Fig.~\ref{fig_6}(b)) and a downsampling operation. The DCABlock incorporates channel attention based on the DCRB (Figure 5). Concerning the local branch, $F_e^{h}$ and $F_d^{h}$ share the same local branch structure. However, considering that optical flow features $F_o^{h}$ contain directional information, the downsampling is substituted with a 2D Directional Max Pooling (2D DMP)\cite{39} (Fig.~\ref{fig_6}(c)). By utilizing the sign information of pixel values, it separates positive and negative directional components, performs max pooling based on directional information, and then recombines them. After processing through the local branch, features $F_L^{4}$ are obtained for each modality. After the DBFA process across all modalities, the global feature $F_G^{4}$ and the local feature $F_L^{4}$ are concatenated to yield the unimodal feature $F_m^{U} = \operatorname{\mathit{concat}}\!\left(F_G^{4}, F_L^{4}\right)$.


\begin{figure}[!t]
\centering
\includegraphics[width=3.5in]{fig6}
\caption{The network architecture diagram of the DBFA block.}
\label{fig_6}
\end{figure}

\subsection{Multimodal Fusion}

To alleviate the problem of modal heterogeneity, a two-stage fusion strategy is employed during the multimodal fusion network. The ADWF and BGIF blocks are devised to accomplish cross-modal learning from initial integration to deep interaction, which tackles the heterogeneity of distinct modalities in both feature space and temporal sequences. This methodology identifies latent complementary features associated with emotion across modalities, thus improving information representation capabilities. 

\subsubsection{The ADWF Block}

The ADWF block serves as the first-stage fusion module for unimodal features $F_o^{U}$, $F_e^{U}$ and $F_d^{U}$, with its structure illustrated in Fig.~\ref{fig_7}. First, the three unimodal features are concatenated to obtain $F_t$:

\begin{equation}
\label{eq_5}
F_t = \operatorname{\mathit{concat}}\!\left(F_o^{U}, F_e^{U}, F_d^{U}\right)
\end{equation}

Subsequently, after normalization, the features are passed through a linear layer, followed by the GELU activation function, dropout, a linear layer and Softmax activation function. So, the weighted scores $w' = [w_1', w_2', w_3']$ for each modality feature are calculated as follows: 

\begin{equation}
\label{eq_6}
\begin{split}
w' =
\operatorname{\mathit{softmax}}\Bigl(&
\operatorname{\mathit{Linear}}\bigl(
\operatorname{\mathit{Dropout}}\bigl(
\operatorname{\mathit{GELU}}\bigl(\\ &
\operatorname{\mathit{Linear}}\bigl(
\operatorname{\mathit{Norm}}(F_t)\bigr) \bigr) \bigr) \bigr) \Bigr)
\end{split}
\end{equation}

Finally, based on $w'$, the unimodal features are combined through a weighted summation to obtain the adaptive fusion feature $F_A$, that is:

\begin{equation}
\label{eq_7}
F_A = w_1' \times F_o^{U} + w_2' \times F_e^{U} + w_3' \times F_d^{U}
\end{equation}


\begin{figure}[!t]
\centering
\includegraphics[width=2.5in]{fig7}
\caption{The network architecture diagram of the ADWF block.}
\label{fig_7}
\end{figure}


\subsubsection{The BGIF Block}

To attain the integration of unimodal features $F_m^{U}$ and multimodal features $F_A$, the BGIF block founded on multi-head attention is constructed, as depicted in Fig.~\ref{fig_8}. Through the utilization of two parallel multi-head attention mechanisms, the integration of $F_m^{U}$ and $F_A$ is accomplished.

\begin{figure}[!t]
\centering
\includegraphics[width=2.5in]{fig8}
\caption{The network architecture diagram of the BGIF Block.}
\label{fig_8}
\end{figure}

Multi-head attention 1 takes the unimodal feature $F_m^{U}$ as the input $Q$, denoted as $Q_U$, while utilizing the multimodal feature $F_A$ as the inputs $K$ and $V$, denoted as $K_A$ and $V_A$. Conversely, multi-head attention 2 adopts the multimodal feature $F_A$ as the input $Q$, denoted as $Q_A$ and the unimodal feature $F_m^{U}$ as the inputs $K$ and $V$, denoted as $K_U$ and $V_U$. Firstly, we calculate the scaled dot-product attention as follows:

\begin{equation}
\label{eq_8}
A_{1}^{i}
=
\operatorname{\mathit{softmax}}\!\left(
\frac{Q_{U} K_{A}^{\mathrm{T}}}{\sqrt{d_{k_A}}}
\right)
V_{A}
\end{equation}

\begin{equation}
\label{eq_9}
A_{2}^{i}
=
\operatorname{\mathit{softmax}}\!\left(
\frac{Q_{A} K_{U}^{\mathrm{T}}}{\sqrt{d_{k_U}}}
\right)
V_{U}
\end{equation}

\noindent where, $i = 1, 2, 3, 4$, $A_{o}^{i}\ (o \in \{1, 2\})$ denotes the feature representation output by the $i\text{-th}$ attention head in multi-head attention $O$. $d_{k_x}\,(x \in \{A, U\})$ denotes the dimension of each attention head.

Then, the outputs of the two multi-head attention layers are concatenated as:

\begin{equation}
\label{eq_10}
\mathrm{\mathit{MHA}_1}
=
\operatorname{\mathit{Concat}}\!\left(
A_{1}^{i}, \ldots, A_{1}^{i}
\right)
W_{1}
\end{equation}


\begin{equation}
\label{eq_11}
\mathrm{\mathit{MHA}_2}
=
\operatorname{\mathit{Concat}}\!\left(
A_{2}^{i}, \ldots, A_{2}^{i}
\right)
W_{2}
\end{equation}

\noindent where, $W_1$ and $W_2$ are the learnable linear mapping matrixes. Next, we perform normalization on $\mathrm{\mathit{MHA}_1}$ and $\mathrm{\mathit{MHA}_2}$ to obtain $F_1$ and $F_2$. Then, they are concatenated to obtain $F_c = \operatorname{\mathit{Concat}}\!\left(F_1, F_2\right)$. Finally, the weight value $g$ are calculated using $g = \operatorname{\mathit{Sigmoid}}\!\left(\operatorname{\mathit{Linear}}\!\left(F_c\right)\right)$, yielding the final output feature $F_{m}^{f}$ by the follow equations:

\begin{equation}
\label{eq_12}
F_s = (1 - g) \cdot F_1 + g \cdot F_2
\end{equation}

\begin{equation}
\label{eq_13}
F_m^{f} = \operatorname{\mathit{Norm}\!\left(\operatorname{\mathit{Linear}}\!\left(F_s\right)\right)}
\end{equation}

Repeat the above process three times until the three unimodal features $F_{o}^{U}$, $F_{e}^{U}$, and $F_{d}^{U}$ sequentially complete bidirectional guidance with the multimodal feature $F_A$, yielding the fused features $F_{o}^{f}$, $F_{e}^{f}$, and $F_{d}^{f}$, respectively. Then, the multimodal feature $F_A$ is concatenated with them, yielding the final fusion feature $F_f = \operatorname{\mathit{concat}}\!\left(F_{o}^{f},F_{e}^{f},F_{d}^{f}, F_A\right)$, which is fed into the classification head to perform emotion category classification.

\subsection{The Classification and Loss Function}

The classification head consists of a fully connected layer, employing the GELU activation function to introduce nonlinearity and incorporating dropout to enhance model stability. Subsequently, the Softmax function is utilized to obtain the probabilities for each expression category. We employ the weighted cross-entropy loss $L_{\mathit{WCE}}$ as the loss function. Assume that the total quantity of the training set is $N$ and the $\mathit{c\text{-th}}$ class comprises $N_c$ samples. Firstly, the class weights $w_c$ can be defined as $w_c = \frac{N}{N_c}$. We allocate relatively higher weights to classes with a smaller number of samples and relatively lower weights to classes with a larger number of samples. This allows the model to accommodate class imbalance and mitigate bias. Secondly, we perform normalization on $w_c$, with the calculation process as follows:

\begin{equation}
\label{eq_14}
\operatorname{\mathit{mean}}(w) = \frac{1}{C}\sum_{c=1}^{C} w_c
\end{equation}

\begin{equation}
\label{eq_15}
\hat{w}_c = \frac{w_c}{\operatorname{\mathit{mean}}(w)}
\end{equation}

\noindent where, $C$ denotes the number of emotion categories, $\hat{w}_c$ denotes the normalized weight, representing the relative importance of each category. Finally, we compute the weighted cross-entropy loss function as follows:

\begin{equation}
\label{eq_16}
L_{\mathit{WCE}}
=
-\frac{1}{N}
\sum_{i=1}^{N}
\sum_{c=1}^{C}
\hat{w}_c \cdot y_i^{c} \\{\mathit{log}\!\left(p_i^{c}\right)}
\end{equation}

\noindent where, $y_i^{c}$ denotes is the true label of the $i\text{-th}$ sample in class $C$, and $p_i^{c}$ represents the predicted label.


\section{Experiments}

\subsection{Datasets}

The proposed method is experimentally validated on the publicly available micro-expression dataset CAS(ME)3, which consists of three parts, including Part A, Part B, and Part C. Part A and Part B are based on the second-generation paradigm of MEs elicitation, providing annotated and unannotated ME samples, respectively. Part C employs the third-generation of ME eliciting paradigm, collecting multimodal data from 31 subjects in a laboratory environment with strictly controlled lighting and background. It labels 165 MEs and 347 MAEs while simultaneously acquiring depth information and collecting multiple types of physiological signals, including EDA, ECG, RSP, PPG, and Voice signals. This provides richer, more ecologically valid data to advance genuine emotion recognition. The experiment utilizes Part C for emotion classification, categorizing emotions into four categories, including negative, positive, surprised, and other, as presented in Table \ref{tab:table1}.

\begin{table}[!t]
\caption{Sample Information of CAS(ME)$^{3}$-PartC\label{tab:table1}}
\centering
\begin{tabular}{cccccc}
\toprule[1.5pt] 
\multirow{2}{*}{\textbf{Dataset}} & \multirow{2}{*}{\textbf{Subject}} & \multirow{2}{*}{\textbf{Modality}} & \multirow{2}{*}{\textbf{ME}} & \multicolumn{2}{c}{\textbf{Emotion}}  \\ 
\cmidrule[1pt]{5-6}
                   \multicolumn{4}{c}{} & \textbf{Class} & \textbf{Counts} \\ 
                   \midrule[1pt] 
\multirow{6}{*}{CAS(ME)$^{3}$-PartC} & \multirow{6}{*}{31} & RGB   & \multirow{6}{*}{165} &  &  \\
                               &                     & Depth &                      & Negative & 99 \\
                               &                     & Voice &                      & Positive & 16 \\
                               &                     & EDA   &                      & Surprise & 30 \\
                               &                     & ECG   &                      & Others & 20 \\
                               &                     & PPG   &                      &   &   \\
\bottomrule[1.5pt]
\end{tabular}
\end{table}

\subsection{Experiments Setup}

All experiments are conducted under the same configuration settings. The operating system is Windows 10, equipped with an Intel i5-10400F processor and an NVIDIA GeForce RTX 3080 GPU. The experimental environment is configured using Python 3.8. During training, the batch size is set to 32, the initial learning rate is set to 0.0001, and the AdamW optimizer is employed for optimization.

To evaluate the model's generalization performance and avoid biases arising from individual subject differences, all experiments employed the Leave-One-Subject-Out (LOSO) cross-validation method. To assess the effects of the few-sample problem and class imbalance in the dataset, we employ the Unweighted F1-score (UF1) and Unweighted Average Recall (UAR) as evaluation metrics, and they are defined as: 

\begin{equation}
\label{eq_17}
UF1 = \frac{1}{C} \sum_{c=1}^{C}
\frac{2 \times TP_c}{2 \times TP_c + FP_c + FN_c}
\end{equation}

\begin{equation}
\label{eq_18}
UAR = \frac{1}{C} \sum_{c=1}^{C} \frac{TP_c}{N_c}
\end{equation}

\noindent where, $TP_c$, $FP_c$, and $FN_c$ represent the true positives, false positives, and false negatives for category $c$, respectively.

\subsection{Comparative Results}

AlexNet\cite{24} was originally introduced as a baseline model in conjunction with the release of the dataset. To evaluate the performance of the proposed method, we compare the TMFIF-Net with AlexNet\cite{24}, STSTNet\cite{18}, MulT\cite{40}, MMGCN\cite{41}, MultiEMO\cite{42}, MITRL\cite{43}, HTNet\cite{23}, CCMA\cite{44}, AEFNet\cite{45}, MMMER\cite{7}, and DMINet\cite{46}. Among these, due to the unavailability of source code for AlexNet, MMMER, and DMINet, the reproduced results fell short of the published outcomes. To ensure a fair comparison, we directly cite the experimental results from the original papers. Furthermore, CAS(ME)$^{3}$ is currently the only multimodal micro-expression dataset that includes both physiological signals and depth information. Therefore, this paper evaluates the model’s performance on this dataset. The results in Table \ref{tab:table2} demonstrate that the AlexNet method benefited from depth information, which can contribute to more effective motion feature extraction for micro-expressions in the RGB modality. The MMMER confirmed that the multimodal model of RGB and physiological signals has achieved significant performance improvements compared to the AlexNet method, with UF1 and UAR increasing by 34.6\% and 28.2\%, respectively, and further optimized the accuracy of emotion recognition. Compared to the AlexNet method, the proposed TMFIF-Net achieves improvements of 41.64\% and 36.73\% in UF1 and UAR, respectively. Relative to the method MMMER approach, UF1 of our method is improved by 7.04\%. The proposed method outperforms DMINet by 5.73\% in terms of UAR, which significantly enhances the robustness of genuine emotion analysis.


\begin{table*}[!t]
\caption{Performance Comparison of Four Classification Models. The optimal and sub-optimal results are bolded and underlined.\label{tab:table2}}
\centering
\begin{tabularx}{\textwidth}{l *{8}{>{\centering\arraybackslash}X}}
\toprule[1.5pt]
\multirow{2}{*}{\textbf{Method}} &
\multicolumn{6}{c}{\textbf{Modality}} &
\multicolumn{2}{c}{\textbf{CAS(ME)$^{3}$-PartC}} \\
\cmidrule[1pt]{2-7} \cmidrule[1pt]{8-9}  
& \textbf{RGB} & \textbf{Depth} & \textbf{Voice} & \textbf{EDA} & \textbf{PPG} & \textbf{ECG}
& \textbf{UF1} & \textbf{UAR} \\ \midrule[1pt]

\multirow{7}{*}{\shortstack[l]{Baseline \cite{24}\\(AlexNet, 2022)}}
          & $\checkmark$ &  &  &  &  &  & 0.2480 & 0.2630 \\
          & $\checkmark$ & $\checkmark$ &  &  &  &  & 0.2960 & 0.2960 \\
          & $\checkmark$ &  & $\checkmark$ &  &  &  & 0.1950 & 0.2160 \\
          & $\checkmark$ &  &  & $\checkmark$ &  &  & 0.2300 & 0.2600 \\
          & $\checkmark$ & $\checkmark$ & $\checkmark$ &  &  &  & 0.2410 & 0.2540 \\
          & $\checkmark$ & $\checkmark$ &  & $\checkmark$ &  &  & 0.2300 & 0.2440 \\
          & $\checkmark$ & $\checkmark$ & $\checkmark$ & $\checkmark$ &  &  & 0.1910 & 0.2230 \\
STSTNet \cite{18} (2019)   & $\checkmark$ & $\checkmark$ &  &  &  & $\checkmark$ & 0.2513 & 0.3791 \\
MulT \cite{40} (2019)   & $\checkmark$ & $\checkmark$ &  &  &  & $\checkmark$ & 0.4089 & 0.3856 \\
MMGCN \cite{41} (2021)   & $\checkmark$ & $\checkmark$ &  &  &  & $\checkmark$ & 0.4503 & 0.4121 \\
MultiEMO \cite{42} (2023)   & $\checkmark$ & $\checkmark$ &  &  &  & $\checkmark$ & 0.4380 & 0.4335 \\
MITRL \cite{43} (2023)   & $\checkmark$ & $\checkmark$ &  &  &  & $\checkmark$ & 0.4744 & 0.4396 \\
HTNet \cite{23} (2024)   & $\checkmark$ & $\checkmark$ &  &  &  & $\checkmark$ & 0.4567 & 0.4295 \\
CCMA \cite{44} (2024)   & $\checkmark$ & $\checkmark$ &  &  &  & $\checkmark$ & 0.3826 & 0.3656 \\
AEFNet \cite{45} (2025)   & $\checkmark$ & $\checkmark$ &  &  &  & $\checkmark$ & 0.2426 & 0.2766 \\
\multirow{7}{*}{MMMER \cite{7} (2025)}
          & $\checkmark$ &  &  &  &  &  & 0.3530 & 0.3450 \\
          & $\checkmark$ & $\checkmark$ &  &  &  &  & 0.3150 & 0.3180 \\
          & $\checkmark$ & $\checkmark$ &  & $\checkmark$ & $\checkmark$ & $\checkmark$ & 0.5860 & 0.5630 \\
          & $\checkmark$ &  &  & $\checkmark$ & $\checkmark$ & $\checkmark$ & \underline{0.6420} & 0.5780 \\
          &  &  &  & $\checkmark$ & $\checkmark$ & $\checkmark$ & 0.5230 & 0.5080 \\

DMINet \cite{46} (2026)   & $\checkmark$ & $\checkmark$ &  &  &  &  & 0.6113 & \underline{0.6060} \\

\textbf{Ours(TMFIF-Net)}  & $\checkmark$ & $\checkmark$ &  &  &  & $\checkmark$ & $\textbf{0.7124}$ & $\textbf{0.6633}$ \\ \bottomrule[1.5pt]
\end{tabularx}
\end{table*}

In this paper, we also compare the computational complexity of the proposed method with HTNet and MMMER. As shown in Table \ref{tab:table3}, the TMFIF-Net method has a markedly lower parameter count than HTNet and MMMER. In terms of FLOPs, the proposed method incurs higher FLOPs than the others, indicating a higher computational cost. Although the parameter count has been minimized, substantial computational resources are still required based on FLOPS requirements. Overall, the TMFIF-Net method improves performance with fewer parameters, but requires more FLOPs, implying a trade-off between model compactness and computational efficiency.

\begin{table}[!t]
\caption{Comparison of model complexity with other methods\label{tab:table3}}
\centering
\setlength{\tabcolsep}{3pt}
\begin{tabular}{l*{4}{c}}
\toprule[1.5pt]
\multirow{1}{*}{\textbf{Method}} & \makecell{\textbf{Parameters}\\\textbf{(M)}}
 & \makecell{\textbf{FLOPS}\\\textbf{(G)}}& \makecell{\textbf{Execution Time} \\\textbf{(ms)}} & \makecell{\textbf{Memory} \\\textbf{(MB)}}\\ 
\midrule[1pt]
HTNet \cite{23}  & 139.50 & 0.21 & 9.57 & 561.39\\
MMMER \cite{7}   & 33.34 & 0.17 & 13.59 & 298.25 \\
\textbf{Ours(TMFIF-Net)}  & 16.26 & 2.24 & 32.00 & 90.17\\
\bottomrule[1.5pt]
\end{tabular}
\end{table}



\subsection{Confusion Matrix Analysis}

The experiment employs a confusion matrix to evaluate and analyze the micro-expression classification task, as illustrated in Fig.~\ref{fig_9}. We employ weighted cross-entropy loss to enhance the recognition accuracy of various emotion categories and mitigate the issue of class imbalance. The negative and other emotion categories demonstrate superior performance. Although the positive and surprise categories exhibit lower accuracy, which is improved by 12\% and 16\%, respectively, compared to using the cross-entropy loss function. There is confusion in each emotion category, but due to the small amount of data and imbalanced class distribution, a proportion of 60\% of the samples belongs to the negative category, so the other three emotion categories tend to be misjudged as negative emotion to a certain extent.

\begin{figure}[!t]
\centering
\includegraphics[width=2.5in]{fig9}
\caption{Confusion Matrix.}
\label{fig_9}
\end{figure}

\subsection{Ablation Study}
\subsubsection{The Contribution Analysis of Different Modality}

To assess the contribution and complementarity of each modality, a series of experiments with diverse input modalities was designed, as presented in Table \ref{tab:table4}. For unimodal inputs, depth and physiological modalities (ECG and PPG) consistently outperform RGB in emotion recognition. Compared to RGB, adopting EDA leads to only a 0.85\% improvement in UF1, indicating that EDA has relatively weaker performance in micro-expression recognition relative to other modalities. In comparison to using RGB as input, the model utilizing depth information as input can attain improvements of 2.74\% and 0.16\% in UF1 and UAR, respectively. Additionally, adopting PPG as input results in a 4.9\% and 1.12\% increase in UF1 and UAR. Moreover, when employing ECG as input, the improvement is enhanced to 7.55\% and 4.7\% in UF1 and UAR. Nevertheless, in comparison to depth information, PPG and ECG, RGB data is the most readily accessible. Therefore, the RGB modality is retained in the experimental setup. 

A comparison between bimodal and unimodal approaches reveals the superiority of bimodal approaches. In the case of bimodal approaches, the integration of depth information into RGB leads to improvements of 9.39\% and 5.94\% in UF1 and UAR, respectively. Adding PPG to RGB improves UF1 by 9.58\% and UAR by 5.89\%. Replacing PPG with ECG yields larger gains, with UF1 increased by 13.18\% and UAR by 13.10\%. Among all bimodal settings, RGB+ECG achieves the best overall performance. The combination results show that depth information, PPG, and ECG contribute strongly to emotion recognition. When synergized with RGB, these modalities further improve recognition accuracy.

\begin{table}[!t]
\caption{The Contribution Analysis Results of Different Modality Input. Unimodal, Bimodal, and Trimodal refer to single-modality, dual-modality, and three-modality. The optimal and sub-optimal results are bolded and underlined. \label{tab:table4}}
\centering
\setlength{\tabcolsep}{3pt}
\begin{tabular}{l *{7}{c}}
\toprule[1.5pt]
\multirow{2}{*}{\textbf{Method}} &
\multicolumn{5}{c}{\textbf{Modality}} &
\multicolumn{2}{c}{\textbf{CAS(ME)$^{3}$-PartC}} \\
\cmidrule[1pt]{2-6}\cmidrule[1pt]{7-8}
& \textbf{RGB} & \textbf{Depth}  & \textbf{EDA} & \textbf{PPG} & \textbf{ECG}
& \textbf{UF1} & \textbf{UAR} \\ \midrule[1pt]

\multirow{5}{*}{\shortstack[c]{Unimodal}}
          & $\checkmark$ &  &  &  &  & 0.5050 & 0.4966 \\
          &  & $\checkmark$ &  &  &  & 0.5324 & 0.4982 \\
          &  &  & $\checkmark$ &  &  & 0.5135 & 0.4840 \\
          &  &  &  & $\checkmark$ &  & 0.5541 & 0.5078 \\
          &  &  &  &  &  $\checkmark$ & 0.5805 & 0.5436 \\ \midrule[1pt]

\multirow{10}{*}{\shortstack[c]{Bimodal}}
          & $\checkmark$ & $\checkmark$ &  &  &  & 0.5989 & 0.5560 \\
          & $\checkmark$ &  & $\checkmark$ &  &  & 0.5302 & 0.5042 \\
          & $\checkmark$ &  &  & $\checkmark$ &  & 0.6008 & 0.5555 \\
          & $\checkmark$ &  &  &  & $\checkmark$ & 
          \underline{0.6368} & \underline{0.6276} \\
          &  & $\checkmark$ & $\checkmark$ &  &   & 0.5855 & 0.5383 \\
          &  & $\checkmark$ &  & $\checkmark$ &   & 0.6283 & 0.5957 \\
          &  & $\checkmark$ &  &  &  $\checkmark$ & 0.6346 & 0.6048 \\
          &  &  & $\checkmark$ & $\checkmark$ &   & 0.5576 & 0.5155 \\
          &  &  & $\checkmark$ &  & $\checkmark$  & 0.6014 & 0.5555 \\
          &  &  &  & $\checkmark$ & $\checkmark$  & 0.6317 & 0.6069 \\
          \midrule[1pt]

\multirow{6}{*}{\shortstack[c]{Trimodal}}
          & $\checkmark$ &  & $\checkmark$ & $\checkmark$ &  & 0.5711 & 0.5502 \\
          & $\checkmark$ &  & $\checkmark$ &  & $\checkmark$ & 0.6027 & 0.5752 \\
          & $\checkmark$ &  &  & $\checkmark$ & $\checkmark$ & 0.6084 & 0.5758 \\
          & $\checkmark$ & $\checkmark$ & $\checkmark$ &  &  & 0.5762 & 0.5441 \\
          & $\checkmark$ & $\checkmark$ &  & $\checkmark$ &  & 0.6256 & 0.5786
          \\
          & $\checkmark$ & $\checkmark$ &  &  & $\checkmark$ & $\textbf{0.7124}$ & $\textbf{0.6633}$ \\
\bottomrule[1.5pt]
\end{tabular}
\end{table}

Considering the overall performance of bimodal fusion, incorporating PPG and ECG into the RGB modality may suppress the information representation of the RGB modality. However, the trimodal fusion approach integrating RGB, depth information, and ECG attains the optimal outcomes across all evaluation metrics and notably outperforms other fusion methods. This phenomenon might be attributed to the substantial individual variations and hysteresis of EDA. Additionally, PPG is readily influenced by pressure and light exposure. As a result, the contributions of EDA and PPG are less significant than those of ECG. 


\subsubsection{The Contribution Analysis of How to Convert ECG into Two-Dimensional Images}

This experiment compared three conversion methods, including GADF\cite{34}, GASF\cite{34}, and RP\cite{47}. The results are shown in Table \ref{tab:table5}. GADF captures the local fluctuation of the ECG by calculating the angle difference between two points in the ECG. GASF reflects the overall change trend by calculating the angle sum between two points in the ECG. RP calculates the recursiveness of time series and generates a distance matrix to reflect the periodic characteristics of the ECG.  Compared with RP, GADF can improve the UF1 and UAR by 9.8\% and 8.2\%, respectively. The findings indicate that the RP computes the similarity among points within the time series, thereby generating a sparse distance matrix. This, in turn, results in a restricted capacity to capture local variations in ECG and constrains its feature representation. Therefore, the GADF based on angle difference is more suitable to capture the crucial feature information of emotional variations and shows superior performance.


\begin{table}[!t]
\caption{Comparison of Conversion Methods. The optimal and sub-optimal results are bolded and underlined.\label{tab:table5}}
\centering
\setlength{\tabcolsep}{27pt}
\begin{tabular}{lcc}
\toprule[1.5pt]
\multirow{2}{*}{\textbf{Method}} &
\multicolumn{2}{c}{\textbf{CAS(ME)$^{3}$-PartC}} \\
\cmidrule[1pt]{2-3}
& \textbf{UF1} & \textbf{UAR} \\
\midrule
GASF  & 0.5559 & 0.5415 \\
RP   & \underline{0.6144} & \underline{0.5810} \\
GADF  & $\textbf{0.7124}$ & $\textbf{0.6633}$ \\
\bottomrule[1.5pt]
\end{tabular}
\end{table}

\subsubsection{Ablation Experiments of Each Module in the Model}

To verify the effectiveness of each module in the proposed method, experiments involving diverse modules are carried out, and the outcomes are presented in Table \ref{tab:table6}. It is evident that the DBFA block contributes significantly to enhancing emotion classification performance, making UF1 achieve the optimal performance, while UAR yields the second-best results among models A-D. This improvement can be attributed to the DBFA block, which exploits parallel branches to extract complementary features. The local branch employs convolutions and channel attention to capture local features, while the global branch effectively addresses long-range dependencies by capturing global contextual information, thereby enhancing the representational capacity of the model. Among models A-D, the BGIF block achieves the best UAR, outperforming the DBFA block by 0.56\%. Analytically, this is likely due to the BGIF block, which promotes balanced multimodal feature integration. Compared to model E, model F achieves improvements of 2.51\% and 3.29\% in UF1 and UAR, respectively, indicating that different modalities possess distinct emotional information. The ADWF block is used to assign the modality-specific weight, which can effectively enhance the robustness of recognition. In comparison with model F, UF1 and UAR of model G are increased by 6.21\% and 5.67\%, respectively. These results demonstrate that the interaction between unimodal features and multimodal features is crucial for genuine emotion recognition.

\begin{table}[!t]
\caption{Module Impact Analysis. The optimal and sub-optimal results are bolded and underlined.\label{tab:table6}}
\centering
\begin{tabular}{l *{6}{c}}
\toprule[1.5pt]
\multirow{2}{*}{\textbf{Model}} & \multirow{2}{*}{\textbf{HSFE}} & \multirow{2}{*}{\textbf{DBFA}} & \multirow{2}{*}{\textbf{ADWF}} & \multirow{2}{*}{\textbf{BGIF}} & \multicolumn{2}{c}{\textbf{CAS(ME)$^{3}$-PartC}} \\
\cmidrule[1pt]{6-7}
\multicolumn{5}{c}{} & \textbf{UF1} & \textbf{UAR} \\ \midrule[1pt]

A   & $\checkmark$ &  &  &  & 0.3784 & 0.3982 \\
B   &  & $\checkmark$ &  &  & 0.5820 & 0.5660 \\
C   &  &  & $\checkmark$ &  & 0.4714 & 0.4639 \\
D   &  &  &  & $\checkmark$ & 0.5215 & 0.5716 \\
E   & $\checkmark$ & $\checkmark$ &  &  & 0.6252 & 0.5737 \\
F   & $\checkmark$ & $\checkmark$ & $\checkmark$ &  & \underline{0.6503} & \underline{0.6066} \\
G   & $\checkmark$ & $\checkmark$ & $\checkmark$ & $\checkmark$ & $\textbf{0.7124}$ & $\textbf{0.6633}$ \\
\bottomrule[1.5pt]
\end{tabular}
\end{table}

\subsubsection{BGIF Block Analysis}

To verify the efficacy of the BGIF block, a comparison is performed with the Unidirectional Guided Interaction Fusion (UGIF) strategy that utilizes a single branch of multi-head attention. The detailed setting information and results are presented in Table \ref{tab:table7}. Results in Table \ref{tab:table7} demonstrate that the BGIF block achieves optimal performance across all metrics. In comparison with UGIF\_F, the BGIF block exhibits significant performance gains, with UF1 and UAR improving by 14.02\% and 13.44\%, respectively. Moreover, the BGIF block achieves improvements of 5.16\% in UF1 and 4.60\% in UAR relative to UGIF\_S. This indicates that bidirectional guidance more effectively captures inter-modal dependencies, strengthens feature complementarity, and enhances the effectiveness of information fusion. Compared to BGIF\_S, the BGIF block can attain improvements of 3.17\% and 2.26\% in UF1 and UAR, respectively, demonstrating that bidirectional guidance between multimodal and unimodal features is superior to unimodal pairwise guidance. This improvement may stem from the bidirectional interaction, facilitating more thorough information exchange. By combining multimodal representations with unimodal strengths in a secondary fusion stage, more discriminative information is preserved, leading to improved emotion recognition performance.

\begin{table}[!t]
\caption{The Results of the BGIF Block Analysis. The optimal and sub-optimal results are bolded and underlined.\label{tab:table7}}
\centering
\setlength{\tabcolsep}{4.7pt}
\begin{tabular}{l *{8}{c}}
\toprule[1.5pt]
\multirow{2}{*}{\textbf{Module}} 
& \multicolumn{3}{c}{ \makecell{\textbf{Multi-Head}\\\textbf{Attention 1}}}
& \multicolumn{3}{c}{ \makecell{\textbf{Multi-Head}\\\textbf{Attention 2}}} 
& \multicolumn{2}{c}{\textbf{CAS(ME)$^{3}$-PartC}} \\
\cmidrule[1pt]{2-4} \cmidrule[1pt]{5-7} \cmidrule[1pt]{8-9}
& \textbf{Q} & \textbf{K} & \textbf{V} 
& \textbf{Q} & \textbf{K} & \textbf{V}
& \textbf{UF1} & \textbf{UAR} \\ 

\midrule[1pt]

UGIF\_F    & $F_m^{U}$ & $F_A$ & $F_A$ &  &  &  & 0.5722 & 0.5289 \\
UGIF\_S    & $F_A$ & $F_m^{U}$ & $F_m^{U}$ &  &  &  & 0.6608 & 0.6173 \\
BGIF\_S    & $F_m^{U}$ & $F_m^{U}$ & $F_m^{U}$ & $F_m^{U}$ & $F_m^{U}$ & $F_m^{U}$ & \underline{0.6807} & \underline{0.6407} \\
BGIF      & $F_m^{U}$ & $F_A$ & $F_A$ & $F_A$ & $F_m^{U}$ & $F_m^{U}$ & $\textbf{0.7124}$ & $\textbf{0.6633}$ \\
\bottomrule[1.5pt]
\end{tabular}
\end{table}


\subsubsection{DBFA Performance Analysis}

\begin{figure}[!t]
\centering
\includegraphics[width=2.5in]{fig10}
\caption{Analysis of unimodal features extracted by dual branches. G denotes the Global branch, L denotes the Local branch.}
\label{fig_10}
\end{figure}

In order to verify the effectiveness of the DBFA block, the same network architecture is used for dual-branch feature extraction for each modality, and the 2D DMP method is not incorporated into the local branch of each modality. The experiment compares single-branch and dual-branch feature extraction to analyze the effect. The experimental results in Fig.~\ref{fig_10} demonstrate that employing a dual-branch network architecture for unimodal deep feature extraction significantly enhances performance metrics. Compared with Model 1, Model 13 achieves improvements of 2.68\% in UF1 and 1.16\% in UAR. In addition, Model 13 exhibits an increase of 6.13\% in UF1 and 3.76\% in UAR compared to Model 2, indicating that the dual-branch network achieves the optimal performance and shows superior effectiveness in unimodal feature extraction compared to single-branch models using G or L. In comparison with Model 4, Model 10 achieves increases of 0.27\% in UF1 and 1.89\% in UAR. Relative to Model 11, Model 5 improves UF1 and UAR by 6.23\% and 6.55\%, respectively. Compared with Model 12, Model 7 yields improvements of 2.93\% and 3.52\% in UF1 and UAR, respectively. These comparisons demonstrate the advantage of using the L branch for depth information and the G branch for ECG and RGB. Model 9 obtains increases of 1.90\% in UF1 and 3.04\% in UAR relative to Model 3, indicating that adopting the L branch for depth information and ECG signals tends to outperform the G branch under the RGB dual-branch architecture. Among all compared methods, Model 8 attains the second-best performance. In summary, it is evident that the parallel dual-branch network architecture provides more complementary information, enhances feature representation, and effectively improves emotion recognition performance.


\subsubsection{2D DMP Module Analysis}

To evaluate the effectiveness of replacing downsampling with 2D DMP when processing optical flow features in the DBFA Block, comparative experiments are conducted on local branches while keeping other modules unchanged. The experimental results are presented in Table \ref{tab:table8}. The experimental results demonstrate that incorporating directional information into the local branches of the RGB effectively enhances the performance metrics of UF1 and UAR, achieving respective improvements of 3.14\% and 3.32\% relative to the baseline without 2D DMP across all modalities. Analysis suggests that this may be attributable to the fact that optical flow features represent the displacement of micro-expression over time, and the signs of adjacent regions with obvious motion characteristics may be opposite, while directional perception can effectively preserve information regarding facial expression changes. For the ECG, directional pooling operations may compromise temporal information, failing to capture periodic variations. However, the depth information reflects the changes of the front and back direction and gradient, and focuses more on the spatial distribution. Therefore, directional differentiation may result in the loss of certain depth information.

\begin{table}[!t]
\caption{2D DMP Module Analysis. The optimal and sub-optimal results are bolded and underlined.\label{tab:table8}}
\centering
\setlength{\tabcolsep}{8pt}
\begin{tabular}{l *{5}{c}}
\toprule[1.5pt]
\multirow{2}{*}{\textbf{Module}} & \multirow{2}{*}{\textbf{RGB}} & \multirow{2}{*}{\textbf{ECG}} & \multirow{2}{*}{\textbf{Depth}} & \multicolumn{2}{c}{\textbf{CAS(ME)$^{3}$-PartC}} \\
\cmidrule[1pt]{5-6}
\multicolumn{4}{c}{} & \textbf{UF1} & \textbf{UAR} \\ \midrule[1pt]

\multirow{4}{*}{\shortstack[c]{2D DMP}}   
   &  &  &  & \underline{0.6810} & \underline{0.6301} \\
   & $\checkmark$ &  &  & $\textbf{0.7124}$ & $\textbf{0.6633}$ \\
   &  & $\checkmark$ &  & 0.6032 & 0.5716 \\
   &  &  & $\checkmark$ & 0.6011 & 0.5455 \\
\bottomrule[1.5pt]
\end{tabular}
\end{table}


\subsubsection{ADWF Module Analysis}

To validate the effectiveness of the ADWF module for unimodal feature fusion, its performance is compared with that of direct concatenation. The results are presented in Table \ref{tab:table9}. Table \ref{tab:table9} experimental results indicate that ADWF outperforms the direct concatenation method, achieving improvements of 3.77\% in UF1 and 3.16\% in UAR, respectively. While the fusion approach involving direct concatenation of unimodal features achieved relatively favorable performance, it failed to account for the varying degrees of importance of each modality. Feature-based weights are learned in ADWF to adjust modality contributions. Shallow interactions across modalities are thus introduced, leading to improved emotion recognition accuracy.

\begin{table}[!t]
\caption{ADWF Module Analysis. The optimal and sub-optimal results are bolded and underlined.\label{tab:table9}}
\centering
\setlength{\tabcolsep}{26pt}
\begin{tabular}{l *{2}{c}}
\toprule[1.5pt]
\multirow{2}{*}{\textbf{Module}} &
\multicolumn{2}{c}{\textbf{CAS(ME)$^{3}$-PartC}} \\
\cmidrule[1pt]{2-3}
\multicolumn{1}{c}{} & \textbf{UF1} & \textbf{UAR} \\ \midrule[1pt]
Concat  & \underline{0.6747} & \underline{0.6317} \\
ADWF   & $\textbf{0.7124}$ & $\textbf{0.6633}$ \\
\bottomrule[1.5pt]
\end{tabular}
\end{table}

\section{Conclusion}

Multimodal emotion recognition combines heterogeneous modalities to enrich feature representations. This capability is important for improving human–computer interaction. To fully leverage multimodal information and capture intrinsic inter-modal correlations, we propose TMFIF-Net, in which visual and physiological emotion features are jointly modeled. Complementary affective cues are captured from RGB, depth, and ECG modalities. A progressive hierarchical feature extraction network is designed, employing the DBFA block to extract the most discriminative features within each modality. Cross-modal variations in latent emotional cues and feature representations motivate the adoption of a two-stage fusion strategy. An ADWF block dynamically balances contributions across modalities through weight adjustment. A BGIF block based on cross-modal fusion via attention mechanisms enables bidirectional understanding between unimodal and multimodal interactions, enhancing MER performance. Although the experiments conducted in this paper have achieved some progress, they have failed to resolve the issue of emotion category imbalance. In future work, the performance of minority categories will be further explored, independent feature extractors will be designed for the specific emotional information of each modality, and other EDA and PPG physiological signals will be used in synergy to improve the emotion recognition accuracy. Furthermore, efforts will be directed towards addressing the issue of missing modality, thereby enhancing the system's adaptive ability in the face of incomplete data to overcome the limitations of real-world emotion analysis in complex scenarios.



%{\appendices
%\section*{Proof of the First Zonklar Equation}
%Appendix one text goes here.
% You can choose not to have a title for an appendix if you want by leaving the argument blank
%\section*{Proof of the Second Zonklar Equation}
%Appendix two text goes here.}


\bibliographystyle{IEEEtran}
\bibliography{refs}

%\bibitem{ref1}








\vfill

\end{document}


